\begin{ouverture}[calcul]{8}{Prendre une fraction d’une quantité}{Partie 2 · Fractions et nombres décimaux}
{Opérateur · Deux tiers de… · Problèmes}
\objectif{Calculer une fraction d’une quantité ou d’une grandeur.}
\objectif{Multiplier un entier par une fraction.}
\objectif{Retrouver le tout à partir d’une part dans des cas simples.}
\prerequis{Fractions et division euclidienne.}
\end{ouverture}
\pagestyle{otMain}\essentiel\fontsize{11}{15}\selectfont
\phantomsection\label{cours-8-1}
\essSection{1}{Une fraction qui agit}
\begin{definition}Prendre $\frac23$ de 12, c’est partager 12 en 3 parts égales, puis prendre 2 parts : $12:3=4$, puis $4\times2=8$. On note aussi ce calcul $\frac23\times12=8$.\end{definition}
\begin{exemple}Trois quarts de 100 mètres : $100:4=25$, puis $25\times3=75$ mètres. Deux tiers de 12 euros : $12:3=4$, puis $4\times2=8$ euros.\end{exemple}
\begin{propriete}Le produit d’un entier par une fraction peut se calculer dans l’ordre le plus simple : $\frac34\times60$ se calcule $60:4=15$ puis $15\times3=45$ ; ou $3\times60=180$ puis $180:4=45$ : même résultat.\end{propriete}
\begin{savoirfaire}[Encadrer le résultat attendu]Contrôler l’ordre de grandeur : $\frac23$ de 12 doit être plus petit que 12 (fraction inférieure à 1) et plus grand que la moitié de 12, car $\frac23>\frac12$.\end{savoirfaire}
\clearpage
\phantomsection\label{cours-8-2}
\essSection{2}{Du tout à la part, et retour}
\begin{savoirfaire}[Remonter à l’unité de part]Retrouver le tout : si $\frac13$ d’une somme vaut 15 euros, le tout vaut $15\times3=45$ euros. Si $\frac25$ d’un trajet font 30 km, un cinquième fait $30:2=15$ km, et le trajet entier $15\times5=75$ km.\end{savoirfaire}
\begin{exemple}Une classe de 28 élèves : les $\frac47$ sont des filles, soit $28:7=4$, puis $4\times4=16$ filles ; il y a donc $28-16=12$ garçons. Le contrôle $16+12=28$ rassure.\end{exemple}
\begin{attention}La fraction agit sur le tout annoncé, pas sur un reste : « $\frac14$ de 60, puis $\frac13$ du reste » fait $15$, puis $\frac13$ de 45, soit 15 encore. Lire à chaque étape sur quoi porte la fraction.\end{attention}
\begin{propriete}Prendre $\frac55$ d’une quantité la laisse entière ; prendre $\frac{10}5$, c’est la doubler : une fraction plus grande que 1 agrandit.\end{propriete}

\clearpage\rubrique{Découverte}
\begin{blocLivre}{ex-8-1}\begin{exercice}[1]{Des parts simples}
Calculer $\frac13$ de 18, puis $\frac14$ de 32.
\end{exercice}\end{blocLivre}
\begin{blocLivre}{ex-8-2}\begin{exercice}[1]{Deux tiers de douze}
Calculer $\frac23$ de 12 euros, puis $\frac34$ de 100 mètres.
\end{exercice}\end{blocLivre}
\begin{blocLivre}{ex-8-3}\begin{exercice}[1]{Deux chemins, un produit}
Calculer $\frac34\times60$ de deux façons : en partageant d’abord, puis en multipliant d’abord.
\end{exercice}\end{blocLivre}
\begin{blocLivre}{ex-8-4}\begin{exercice}[1]{Encadrer avant de calculer}
Sans calculer exactement, dire entre quels nombres se trouve $\frac23$ de 12, puis vérifier par le calcul.
\end{exercice}\end{blocLivre}
\begin{blocLivre}{ex-8-5}\begin{exercice}[1]{Des fractions d’heure}
Convertir en minutes : $\frac14$ d’heure ; $\frac34$ d’heure ; $\frac1{6}$ d’heure.
\end{exercice}\end{blocLivre}

\rubrique{Entraînement}
\begin{blocLivre}{ex-8-6}\begin{exercice}[2]{La jauge du réservoir}
Un réservoir de 48 litres est rempli aux $\frac38$.\begin{questions}
\item Quel volume d’essence contient-il ?
\item Quel volume manque-t-il pour le remplir ?
\end{questions}
\end{exercice}\end{blocLivre}
\begin{blocLivre}{ex-8-7}\begin{exercice}[2]{Retrouver le tout}
Le tiers d’une somme vaut 15 euros. Quelle est la somme ? Les $\frac25$ d’un trajet font 30 km. Quelle est la longueur du trajet ?
\end{exercice}\end{blocLivre}
\begin{blocLivre}{ex-8-8}\begin{exercice}[2]{La classe de 28}
Dans une classe de 28 élèves, les $\frac47$ sont des filles. Combien y a-t-il de filles ? De garçons ?
\end{exercice}\end{blocLivre}
\begin{blocLivre}{ex-8-9}\begin{exercice}[2]{Plus grand que le tout ?}
Calculer $\frac55$ de 35, puis $\frac{10}5$ de 35. Que remarque-t-on ?
\end{exercice}\end{blocLivre}
\begin{blocLivre}{ex-8-10}\begin{exercice}[2]{L’étape de la fraction}
Sur 60 euros, Tom dépense $\frac14$ le matin, puis $\frac13$ du reste l’après-midi. Combien dépense-t-il chaque fois, et que lui reste-t-il ?
\end{exercice}\end{blocLivre}

\rubrique{Problèmes}
\begin{blocLivre}{ex-8-11}\begin{exercice}[3]{Le jardin potager}
Un jardin de 120 m² consacre $\frac25$ de sa surface aux légumes et $\frac3{10}$ aux fleurs. Quelle surface occupe chaque partie ? Quelle surface reste-t-il en pelouse ?
\end{exercice}\end{blocLivre}
\begin{blocLivre}{ex-8-12}\begin{exercice}[3]{Le prix soldé}
Un blouson coûte 60 euros. Pendant les soldes, il est vendu aux $\frac34$ de son prix. Quel est le nouveau prix ? Quelle économie cela fait-il ?
\end{exercice}\end{blocLivre}
\begin{blocLivre}{ex-8-13}\begin{exercice}[3]{La collecte}
Une école veut collecter 600 livres. Elle en a déjà reçu les $\frac56$. Combien de livres manquent-ils ?
\end{exercice}\end{blocLivre}
\begin{blocLivre}{ex-8-14}\begin{exercice}[3]{Qui a raison ?}
Pour calculer $\frac23$ de 18, Aya calcule $18:3\times2$ et Noé calcule $18\times2:3$. Qui a raison ? Justifier en effectuant les deux calculs.
\end{exercice}\end{blocLivre}
\begin{blocLivre}{ex-8-15}\begin{exercice}[3]{Le tout mystère}
Après avoir dépensé les $\frac34$ de ses économies, Jade possède encore 12 euros. Quel était le montant de ses économies ?
\end{exercice}\end{blocLivre}

\clearpage
